\begin{textcolorbox}[Prompt for Logical Gap Judge]
\textbf{Task:} Evaluate the logical rigor of a proposed solution to a mathematical inequality problem. Focus on whether the reasoning includes non-trivial claims made without justification, logical leaps, or unsupported assertions, while allowing for valid optimization-based, algebraic, or analytic analysis when properly demonstrated.
This prompt does not evaluate whether the direction of an inequality was justified using toy cases, special values, or asymptotic behavior; that aspect is handled separately. \\

\textbf{Instructions:}

1. Carefully read the entire reasoning process used to solve the inequality.

2. Identify whether the solution includes:

\quad - Any non-obvious (non-trivial) claims or transformations without justification or explanation.

\quad - Any logical gaps or skipped steps that lead to intermediate or final conclusions.\\
3. All significant transformations—especially involving inequalities, bounds, or extremal behavior—must be supported by: algebraic manipulation, well-known identities or theorems, valid analytical tools (e.g., convexity, derivatives, limits) or step-by-step numeric or symbolic reasoning.

4. Optimization methods (e.g., Lagrange multipliers, derivative-based analysis) are valid only if the analysis is explicitly shown:

\quad - If a solution invokes optimization or analytical techniques, it must demonstrate key steps, derivative conditions, or critical point verification.

\quad - Statements such as ``solving the constrained optimization problem confirms...'' without any derivation or argument are considered unjustified.

\quad - You do not need to assess whether toy cases, special values, or extreme behavior were used to infer the inequality direction. That responsibility lies outside the scope of this Judge.\\
5. Simple algebra or widely known transformations (e.g., AM-GM, factoring identities) may be used without full derivation.\\
6. The goal is to assess whether each important conclusion within the reasoning—not just the final answer—is logically supported and rigorously justified.\\

\textbf{Output Format:}

\texttt{<Analysis>:} Step-by-step explanation of whether the reasoning is logically sound. Highlight any unjustified claims or skipped steps, unless they are supported by valid asymptotic, numeric, or analytic reasoning. \\
\texttt{<Flagged Reasoning Step (if applicable)>:} Quote or summarize the specific step(s) where an unjustified claim or logical leap occurred. \\
\texttt{<Answer>:} \texttt{True} or \texttt{False}. \texttt{True} if the reasoning is valid; \texttt{False} if it contains unjustified steps or unsupported claims. \\

\textbf{Key Considerations:}

1. Focus on whether each major step (not just the final answer) is logically justified.\\
2. Non-trivial algebraic identities or inequality steps must be explained unless they are well-known.\\
3. Minor simplifications and standard techniques are acceptable without proof.\\
4. Do not flag steps involving toy cases, extreme values, or special substitutions used to infer inequality direction—those are out of scope here.\\
5. Claims like ``a numerical check shows'' must include specific values, results, or graphs to be valid.\\
6. Optimization-based arguments (e.g., Lagrange multipliers, critical point methods) must include demonstrated steps or analytic structure. If only the method is named but not applied, the reasoning should be flagged.\\
7. Do not flag steps used solely for equality verification, sharpness testing, or illustration.\\
8. Do not provide improvement suggestions—simply judge whether the logic is valid or flawed as presented.\\

\textbf{Examples of Inputs and Outputs:}

\texttt{\{examples\}}\\

\textbf{Now analyze the following problem and solution:}

Original Problem:
\texttt{\{query\}}

Solution:
\texttt{\{response\}}

\end{textcolorbox}


